%!TEX root = ./../../thesis-trier.tex

Next, we investigate how $A$-primitivity relates to transitivity
and, more generally, possible decompositions of the permutation domain
into orbits. Lastly, we introduce strongly $\{a,b\}$-primitive permutation groups.

\begin{proposition}
	\label{sync_max:prop:ab_non_congruent_then_two_orbits}
	Let $G \le_G \mathcal S_n$
	and	$A \subseteq [n]$ with $|A| \ge 2$.
	If $G$ is $A$-primitive, then $[n] = \bigcup_{a \in A} \Orbit{a}{G}$.
\end{proposition}
\begin{proof}
	The set $\bigcup_{a \in A} \Orbit{a}{G}$ forms a block containing
	$A$. Hence it equals $[n]$. 
\end{proof}

The next example illustrates \cref{sync_max:prop:ab_non_congruent_then_two_orbits} 
with an intransitive group that is $A$-primitive
for a $2$-set $A$.

\begin{example}\label{sync_max:ex:abprim_intransitive}
	Let $G = \langle \{ g, h\} \rangle_G$
	where $g = (0 ~ 1 ~ 2 ~ 3)$ and $h = (4 ~ 5 ~ 6~ 7)$.
	This group is intransitive. Suppose $\Delta \subseteq [n]$
	is a block with $\{0,4\} \subseteq \Delta$.
	As $4$ is fixed by~$g$, we have $4 \in \applyfunc{\Delta}{g}$ 
	and so we get $\applyfunc{\Delta}{g} = \Delta$.
	Furthermore, we have $i \in \applyfunc{\Delta}{g^i} = \Delta$ for each $i \in \{1,2,3\}$
	so that $\{0,1,2,3\} \subseteq \Delta$.
	Similarly, we find $\{4,5,6,7\} \subseteq \Delta$.
	Hence $\Delta = [n]$ and $G$ is $\{0,4\}$-primitive.
\end{example}

The condition $[n] = \bigcup_{a \in A} \Orbit{a}{G}$
does not imply $A$-primitivity.
For example, if $G \le_G \mathcal S_n$
is transitive and imprimitive,
then for each $a \in [n]$ we have $[n] = \Orbit{a}{G}$, but, by \cref{prim-a-primitive-trans},
it is not $\{a\}$-primitive.
The next example gives an intransitive group
showing that the converse of \cref{sync_max:prop:ab_non_congruent_then_two_orbits}
does not hold in general.

\begin{example}
	Let $G = \langle \{  (0 ~ 1 ~ 2 ~ 3) (4 ~ 5 ~ 6 ~ 7) \} \rangle_G$.
	Then $[8] = \Orbit{0}{G} \cup \Orbit{4}{G}$,
	$\Orbit{0}{G} \ne \Orbit{4}{G}$
	and $\{0,4\}$ is a non-trivial block
	so that $G$ is not $\{0,4\}$-primitive.
\end{example}

Next, we investigate sufficient conditions for transitivity
in terms of $A$-primitivity.

\begin{proposition}
	\label{sync_max:prop:three_non_congruence_implies_transitive} 
	Let $G \le \mathcal S_n$ with $n \ge 3$ and $k \ge 2$.
	Suppose $G$ is $B$-primitive
	for each $k$-subset $B \subseteq A$
	for some $A \subseteq [n]$ with $|A| \ge k(k-1) + 1$.
	Then $G$ is transitive.
\end{proposition}
\begin{proof}
	Let $B \subseteq A$ be some $k$-subset.
	As $G$ is $B$-primitive, by \cref{sync_max:prop:ab_non_congruent_then_two_orbits}, 
	we have $[n] = \bigcup_{b \in B} \Orbit{b}{G}$. 
	By the pigeonhole principle,
	considering the points from $A$
	distributed among the orbits $\Orbit{b}{G}$ for $b \in B$,
	there exists $b \in B$
	such that $| \Orbit{b}{G} \cap A | \ge k$.
	Thus we find a $k$-set $C$ in $\Orbit{b}{G} \cap A$.
	By assumption, $G$ is $C$-primitive.
	Since the orbit $\Orbit{b}{G}$ 
	is a block containing~$C$, we get $\Orbit{b}{G} = [n]$.
\end{proof}

\begin{corollary}
	\label{sync_max:cor:three_non_congruence_implies_transitive} 
	Let $a,b,c \in [n]$ be distinct	and $G \le_G \mathcal S_n$ be $\{a,b\}$-, $\{a,c\}$- and $\{b,c\}$-primitive.
	Then $G$ is transitive.
\end{corollary} 
\begin{proof}
	This follows from~\cref{sync_max:prop:three_non_congruence_implies_transitive} with
	$A = \{a,b,c\}$ and $k = 2$.\footnote{Alternatively,
		from \cref{sync_max:prop:ab_non_congruent_then_two_orbits}
		we get $[n] = \Orbit{a}{G} \cup \Orbit{b}{G} = \Orbit{a}{G} \cup \Orbit{c}{G} = \Orbit{b}{G} \cup \Orbit{c}{G}$.
		Thus, we have $a \in \Orbit{b}{G}$ or $a \in \Orbit{c}{G}$,
		and each case yields $[n] = \Orbit{a}{G}$.}		
\end{proof} 

The next statements connect $A$-primitivity to transitivity.

\begin{proposition} \label{sync_max:prop:ab-prim-transitive}
	Let $G \le_G \mathcal S_n$ and $A \subseteq [n]$ with $|A| \ge 2$.
	Suppose $G$ is $A$-primitive.
	Then~$G$ is transitive if and only if $A \subseteq \Orbit{b}{G}$
	for some $b \in [n]$.
\end{proposition}
\begin{proof}
	First, suppose~$G$ is transitive.
	Pick some $b \in [n]$. Then $\Orbit{b}{G} = [n]$
	and so $A \subseteq \Orbit{b}{G}$.
	
	Next, suppose $A \subseteq \Orbit{b}{G}$.
	Since~$G$ is $A$-primitive,
	by \cref{sync_max:prop:ab_non_congruent_then_two_orbits},
	we have $[n] = \bigcup_{a \in A} \Orbit{a}{G}$.
	Now, $a \in \Orbit{b}{G}$
	implies $\Orbit{a}{G} = \Orbit{b}{G}$.
	Thus, we obtain $[n] = \Orbit{b}{G}$
	and so~$G$ is transitive.	 	
\end{proof}

\begin{corollary}\label{sync_max:prop:ab-prim-transitive}
	Let $G \le_G \mathcal S_n$ be $\{a,b\}$-primitive for distinct $a,b \in [n]$.
	Then there exists $g \in G$
	such that $\applyfunc{a}{g} = b$ if and only if $G$
	is transitive.
\end{corollary}
\begin{proof}
	This follows from \cref{sync_max:prop:ab-prim-transitive}
	since $\{a,b\} \subseteq \Orbit{c}{G}$ for some $c \in [n]$
	is equivalent to the existence of some $g \in G$ with $g(a) = b$.
\end{proof}

\IN{strongly $\{a,b\}$-primitive}
For a $2$-set $\{a,b\}$,
we call $G \le_G \mathcal S_n$ \emph{strongly $\{a,b\}$-primitive}
if $G$ is $\{a,b\}$-primitive and there exists
$g \in G$ such that $\applyfunc{a}{g} = b$ (or equivalently $G$
is $\{a,b\}$-primitive and transitive).
As 
$\applyfunc{b}{g^{-1}} = a$, $G$ is strongly $\{a,b\}$-primitive
if and only if it is strongly $\{b,a\}$-primitive,
so the set notation ``$\{a,b\}$''
is justified here.
%The situation that $A$ is a $2$-set
%and there exists an element mapping one element
%to the other is the situation of interest to us.

\begin{example}
	The permutation group
	from Example~\ref{sync_aut:permgrp:ex:Aprim} 
	is strongly $\{a,b\}$-primitive
	for all distinct $a,b \in [8]$
	with $a \not\equiv b \pmod{2}$.
	The permutation group from Example~\ref{syncaut:permgrp:ex:threesets}
	is strongly $\{ (i,j), (i,j') \}$-primitive
	for $(i,j), (i',j') \in [2] \times [n]$
	with $j \ne j'$.
\end{example} 

As a primitive permutation group over a permutation domain with at least three
points is transitive (\PROP~\ref{sync_max:prop:prim_implies_transitive}), by \PROP~\ref{sync_max:prop:primitive_Aprimitive}
we get the following statement.

\begin{proposition}
	\label{sync_max:prop:prim_stronglyabprim}
	Let $G \le \mathcal S_n$ with $n \ge 3$.
	Then $G$ is primitive if and only if $G$ is strongly $\{a,b\}$-primitive
	for all distinct $a,b \in [n]$.
\end{proposition}
